\SetPageTheme{guided}
\PageHeader
  {Evaluare si recapitulare}
  {Verificam daca poti urmari starea initiala, transformarile si starea finala fara sa pierzi firul.}
  {Recapitulare 1 din 2}

\begin{missionbox}[Scenariu complet]
Pornire:
\begin{lstlisting}[style=codequest]
scor <- 4
vieti <- 3
chei <- 0
\end{lstlisting}
Evenimente:
\begin{enumerate}
\item castigi 6 puncte;
\item pierzi o viata;
\item gasesti 2 chei;
\item folosesti o cheie;
\item mai castigi 3 puncte.
\end{enumerate}
\end{missionbox}

\begin{guidebox}[Rezolvare ghidata]
Completeaza evolutia:
\begin{center}
\begin{tabular}{llll}
\toprule
Pas & scor & vieti & chei \\
\midrule
Initial & 4 & 3 & 0 \\
1 & & & \\
2 & & & \\
3 & & & \\
4 & & & \\
5 & & & \\
\bottomrule
\end{tabular}
\end{center}
\AnswerSpace{5}
\end{guidebox}

\clearpage

\SetPageTheme{challenge}
\PageHeader
  {Autoevaluare}
  {Aici verificam intelegerea, nu viteza. Scrie clar si argumentat.}
  {Recapitulare 2 din 2}

\begin{solobox}[Fisa punctata]
\begin{enumerate}
\item Defineste in propozitie scurta ce este o variabila. \hfill 2p
\item Da doua exemple de operatii valide pe scor. \hfill 2p
\item Construieste o secventa care duce \texttt{energie} de la 9 la 4. \hfill 2p
\item Explica de ce \texttt{x + 1} nu poate inlocui \texttt{x <- x + 1}. \hfill 2p
\item Scrie un mini-scenariu de joc in care folosesti \texttt{scor}, \texttt{vieti}, \texttt{chei}. \hfill 2p
\end{enumerate}
\AnswerSpace{10}
\end{solobox}

\begin{conceptbox}[Interpretarea scorului]
9--10 puncte: fundatia este stabila.

7--8 puncte: mai exerseaza urmarirea starii dupa fiecare pas.

Sub 7 puncte: reia exemplele ghidate si explica cu voce tare fiecare transformare.
\end{conceptbox}

\clearpage
